\chapter*{Sommario}
Sebbene, nei primi anni dell'istruzione secondaria, le ragazze conseguano risultati scolastici in matematica comparabili a quelli dei loro coetanei maschi, persistono significative disparità in termini di autoefficacia matematica, livelli di ansia e orientamento verso future carriere nelle discipline STEM. La presente tesi si inserisce in questo quadro, indagando il potenziale di un seminario di carattere narrativo volto a mettere in discussione la tradizionale rappresentazione della matematica come disciplina prevalentemente maschile.

Il lavoro integra una ricostruzione storica dei contributi offerti dalle donne alla matematica — dall'antichità fino al Novecento, con particolare attenzione alla figura di Maryam Mirzakhani — con uno studio empirico sul campo finalizzato a valutare gli effetti a breve termine di un intervento realizzato in ambito scolastico.

Nel dicembre 2025 è stato realizzato un intervento con disegno pre-test/post-test a gruppo unico presso una scuola secondaria di primo grado di Trieste. I questionari standardizzati somministrati prima e dopo l'intervento sono stati analizzati mediante il software Jamovi, al fine di rilevare eventuali variazioni a breve termine nell'ansia matematica, nell'autoefficacia in matematica, negli atteggiamenti nei confronti della disciplina e nell'adesione esplicita agli stereotipi di genere.

I risultati mostrano che l'intervento è stato associato a una riduzione degli stereotipi di genere espliciti e dei livelli di ansia matematica, nonché a un lieve miglioramento degli atteggiamenti generali degli studenti verso la matematica. L'autoefficacia matematica, invece, è rimasta sostanzialmente stabile, suggerendo che le convinzioni relative alla propria competenza richiedano, per modificarsi, esperienze di padronanza più prolungate e continuative nel tempo.

Nel complesso, lo studio evidenzia come l'integrazione di modelli femminili storici nell'insegnamento della matematica possa costituire una strategia promettente per decostruire la rappresentazione tradizionale della disciplina e promuovere una visione più inclusiva della matematica, intesa come attività umana accessibile, creativa e condivisa.